\begin{proposition}\hypertarget{tag-01Q8}{}\label{fam:src:B.1.2}\label{fam:st:B.1.2}
\StatementRecord{fam:src:B.1.2}{proved}{01Q8}{the definitions and constructions stated in the proof}

The exact set of nonzero parameters for which every ordinary fibre and both normalized finite root fibres of Engel's period-reference quotient are compact complex two-tori is
\[
 0<|u|<e^{-2\pi M}.
\]
This is the parameter domain for the compact family with both smooth finite root models. The condition on ordinary fibres alone at equality is determined by whether a maximizing point of $D_{\rm part}$ is ordinary; this boundary issue is treated in Part C. For an independent reference, the radius is $|\lambda_0|e^{-2\pi M}$.
\end{proposition}


\begin{proof}
\begin{enumerate}[label=\arabic*.,leftmargin=*]
\item The function $N(\xi,W)=|W|\exp(2\pi\operatorname{Im}\mu\,\operatorname{Im}\xi/\operatorname{Im}\tau)$ is invariant under both line-bundle period identifications. The inverse lift multiplies it by $e^{2\pi D_c}$. More directly the real determinant of the four period columns is a nonzero sign times $(\operatorname{Im}\tau)D_c$. Thus the four columns generate a full lattice exactly when $D_c\ne0$, and then the quotient is a compact torus by a closed fundamental parallelepiped. At $D_c=0$ their real span has dimension three; a rank-four lattice in $\mathbb R^4$ cannot result, so the quotient cannot be a compact complex two-torus. The function $D_{\rm part}$ descends continuously to the connected base with the cusp removed and tends to $-\infty$ at the cusp. Its image is therefore $(-\infty,M]$: the maximum is attained, and the intermediate value theorem fills all smaller values. Consequently $D_c$ is nonzero at every base point exactly when $\operatorname{Im}c<-M$. Since $|u|=e^{2\pi\operatorname{Im}c}$, the assertion follows, with the indicated reference scaling.
\end{enumerate}
\end{proof}